\section{Постановка задачи и модель расхождений}
\label{sec:model}

\subsection{Слои сети и измерение расхождений}

Вычисления сети для фиксированного выполнения представим направленным
ациклическим графом (Directed Acyclic Graph, DAG): узлы задают слои
или операции, рёбра~--- передачу тензоров между ними.
Один узел может получать несколько тензоров, а его выход~---
использоваться несколькими узлами (рис.~\ref{fig:transport}).
Цепочка является частным случаем такого графа.

Пронумеруем $L\ge1$ вычислительных узлов в топологическом порядке:
каждый вход поступает от узла с меньшим номером.
Внешние входы объединим в фиксированный вектор
$\Vec{y}^{(0)}\in\mathbb{R}^{n_0}$.
Для узла $l=1,\ldots,L$ множество $\mathcal{P}_l\subseteq\{0,\ldots,l-1\}$
содержит номера непосредственных источников его входов.
Вектор $\Vec{x}^{(l)}\in\mathbb{R}^{\sum_{j\in\mathcal{P}_l}n_j}$
составлен из всех $\Vec{y}^{(j)}$, $j\in\mathcal{P}_l$,
в фиксированном порядке.
Все тензоры развёрнуты в векторы; несколько выходов узла также объединены.
При $n_l>0$ выходных координатах преобразование узла $f^{(l)}$
с параметрами $\Vec{w}^{(l)}\in\mathbb{R}^{p_l}$, $p_l\ge0$, задаёт
\[
    \Vec{y}^{(l)}=f^{(l)}(\Vec{x}^{(l)};\Vec{w}^{(l)})
    \in\mathbb{R}^{n_l}.
\]
Наблюдаемый выход обозначим $\Vec{y}^{(L)}$; если выходов сети несколько,
добавляется конечный узел их объединения без собственной погрешности.
Параметры при сравнении реализаций не меняются.

\begin{figure}[htbp]
    \centering
    \begin{tikzpicture}[
        font=\normalsize,
        box/.style={draw=black!65, rounded corners=2pt,
            fill=babyblue!20, minimum width=16mm, minimum height=8mm},
        arr/.style={-{Stealth[length=4pt]}, line width=0.6pt}
    ]
        \node (in) at (-1.8,0) {$\Vec{y}^{(0)}$};
        \node[box] (f1) at (0,0) {$f^{(1)}$};
        \node[box] (f2) at (3,0.9) {$f^{(2)}$};
        \node[box] (f3) at (3,-0.9) {$f^{(3)}$};
        \node[box] (f4) at (6,0) {$f^{(4)}$};
        \node (out) at (7.8,0) {$\Vec{y}^{(4)}$};
        \draw[arr] (in) -- (f1);
        \draw[arr] (f1.east) -- node[above, sloped] {$\Vec{y}^{(1)}$} (f2.west);
        \draw[arr] (f1.east) -- node[below, sloped] {$\Vec{y}^{(1)}$} (f3.west);
        \draw[arr] (f2.east) -- node[above, sloped] {$\Vec{y}^{(2)}$} (f4.west);
        \draw[arr] (f3.east) -- node[below, sloped] {$\Vec{y}^{(3)}$} (f4.west);
        \draw[arr] (f4) -- (out);
    \end{tikzpicture}
    \caption{Пример графа при $L=4$.
        Выход первого узла используется двумя ветвями;
        четвёртый узел получает два входных тензора.
        Стрелки задают зависимости, а не обязательное параллельное исполнение.}
    \label{fig:transport}
\end{figure}

Слой может быть матричным умножением или содержать его среди других
операций. В слое поэлементной активации матричного умножения нет.
Для матричного умножения обозначим входные матрицы через $A$ и $B$,
результат~--- через $C$, положительные целые размерности~--- $M$, $N$, $K$.
Величина $K$ задаёт число произведений, суммируемых для одного элемента $C$:
\begin{equation}
    \label{eq:gemm}
    C=AB,\qquad
    A\in\mathbb{R}^{M\times K},\quad
    B\in\mathbb{R}^{K\times N},\quad
    C\in\mathbb{R}^{M\times N}.
\end{equation}
Входы и веса заданы в 32-разрядном формате с плавающей точкой
(\texttt{FP32}). К~\eqref{eq:gemm} сводятся линейные проекции,
матричные произведения в механизме внимания и свёртки после развёртки окон
в матрицу (im2col)~\cite{Chellapilla2006}.

Чтобы измерить изменение вычислений, в каждой паре выбирается опорная
реализация. Для сравнения полных путей это PyTorch на CPU без XLA
(далее \emph{CPU-эталон}). Такая опора воспроизводима, но не является
точным вещественным решением или моделью арифметики моделируемой
вычислительной системы.
Выход слоя в опорной реализации обозначим $\Vec{y}^{(l)}$,
в сравниваемой~--- $\widetilde{\Vec{y}}^{(l)}$.
Их разность называется расхождением, или ошибкой относительно опоры.

Разность выходов при одинаковых входах характеризует \emph{локальную}
ошибку узла. Изменённый выход поступает во все использующие его узлы;
каждый преобразует отклонения и может добавить собственную ошибку.
Передачу и преобразование отклонений по графу назовём
\emph{распространением ошибок}.

Для сравнения величины расхождения с масштабом опорного выхода
введём относительную метрику $d_{\mathrm{rel}}^{(l)}$.
Обозначим через $\lVert\cdot\rVert_F$ норму Фробениуса,
то есть корень из суммы квадратов элементов тензора,
а через $\varepsilon_0>0$~--- добавку, предотвращающую деление на ноль:
\begin{equation}
    \label{eq:drel}
    d_{\mathrm{rel}}^{(l)}=
    \frac{\|\widetilde{\Vec{y}}^{(l)}-\Vec{y}^{(l)}\|_F}
         {\|\Vec{y}^{(l)}\|_F+\varepsilon_0},\qquad
    \varepsilon_0=10^{-8}.
\end{equation}
Для вектора эта норма совпадает с евклидовой $\lVert\cdot\rVert_2$.
Добавка $\varepsilon_0$ отличается от машинного эпсилона
$\varepsilon_{\texttt{FP32}}=2^{-23}$.
Нулевое расхождение означает совпадение двух реализаций
в данном сравнении, а не точность относительно вещественной арифметики.

\subsection{Распространение ошибок в прямом проходе}

Для линеаризации $f^{(l)}$~--- дифференцируемая вещественная модель
опорного преобразования, совпадающая с ним на опорном входе;
округление \texttt{FP32} не дифференцируется.
Пусть $\widetilde{\Vec{x}}^{(l)}$ составлен из изменённых выходов
предшественников. Локальное возмущение $\delta_l\in\mathbb{R}^{n_l}$
определим как разность сравниваемого выхода и значения опорной модели
на тех же входах. По определению
\[
    \delta_l:=\widetilde{\Vec{y}}^{(l)}
        -f^{(l)}(\widetilde{\Vec{x}}^{(l)};\Vec{w}^{(l)})
    \quad\Longleftrightarrow\quad
    \widetilde{\Vec{y}}^{(l)}
        =f^{(l)}(\widetilde{\Vec{x}}^{(l)};\Vec{w}^{(l)})+\delta_l.
\]
Это точная запись разности; ошибка может зависеть от входа и параметров.
Для каждого ребра $j\to l$ чувствительность к соответствующему входу
описывает \emph{якобиан}~--- матрица частных производных
$J^{(l,j)}=\partial f^{(l)}/\partial\Vec{y}^{(j)}
\in\mathbb{R}^{n_l\times n_j}$ на опорных тензорах.
Остальные аргументы при этом фиксированы.
Его умножение на изменение входа даёт линейный вклад этого ребра;
вклады всех входов складываются по правилу дифференцирования
сложной функции~\cite[разд.~2]{Petersen2012}.

Основание локального линейного описания~--- покомпонентное разложение
Тейлора по входам~\cite[разд.~10.1.2]{Petersen2012}.
Положим $u_l=\widetilde{\Vec{x}}^{(l)}-\Vec{x}^{(l)}$ и
$\Delta\Vec{y}^{(j)}=\widetilde{\Vec{y}}^{(j)}-\Vec{y}^{(j)}$.
При дифференцируемости опорной модели изменение выхода имеет вид
\[
    \Delta\Vec{y}^{(l)}
    =\sum_{j\in\mathcal{P}_l}J^{(l,j)}\Delta\Vec{y}^{(j)}
       +\delta_l+r_l,\qquad
    \|r_l\|_2=o(\|u_l\|_2).
\]
Здесь $r_l\in\mathbb{R}^{n_l}$~--- остаток линеаризации;
последняя запись означает, что $\|r_l\|_2/\|u_l\|_2\to0$
при $\|u_l\|_2\to0$. Если вторые производные непрерывны и ограничены
в окрестности входа, остаток имеет порядок $O(\|u_l\|_2^2)$.

Эта модель охватывает основные операции сети.
Линейные проекции и свёртки при фиксированных весах дают $r_l=0$.
Если оба аргумента $A$, $B$ в~\eqref{eq:gemm} изменяются
на $\Delta A$, $\Delta B$, остаток равен $\Delta A\Delta B$
и имеет второй порядок по изменениям входов.
Для гладких функций активации и нормализаций, добавляющих положительную
константу к дисперсии или среднему квадрату перед извлечением корня,
применимо разложение Тейлора.
ReLU, то есть $t\mapsto\max(0,t)$, точно линейна по возмущению,
пока не меняются знаки входных координат.
Таким образом, локальное линейное описание существует для этих операций
почти всюду; у ReLU исключение составляют точки
с нулевыми входными координатами.
Это свойство функций, а не гарантия для любой конечной погрешности:
важны величина возмущения, изменение якобиана и расстояние
до границы смены знаков у ReLU.

Определим множество потребителей $\mathcal{S}_j=\{k:j\in\mathcal{P}_k\}$.
Матрица $M_l\in\mathbb{R}^{n_L\times n_l}$ учитывает все пути
от выхода узла $l$ к наблюдаемому выходу, а $I_n$~--- единичная
матрица порядка $n$. Обозначим итоговое отклонение через
$\Delta\Vec{y}^{(L)}=\widetilde{\Vec{y}}^{(L)}-\Vec{y}^{(L)}$,
его линейную часть~--- через $d$, нелинейный остаток~--- через $r_{\mathrm{lin}}$.
Обход графа от выхода к входам даёт
\begin{equation}
    \label{eq:accum}
    \begin{aligned}
        \Delta\Vec{y}^{(L)}&=d+r_{\mathrm{lin}},\qquad
        d=\sum_{l=1}^{L}M_l\delta_l,\\
        M_L&=I_{n_L},\qquad
        M_l=\sum_{k\in\mathcal{S}_l}M_kJ^{(k,l)},\quad 1\le l<L.
    \end{aligned}
\end{equation}
Здесь $r_{\mathrm{lin}}=\sum_l M_l r_l$: остатки также проходят
через последующие операции. Линейное описание пригодно, если норма этой суммы
ниже заданного допуска анализа; малость локальных относительных ошибок
сама по себе этого не доказывает.
Пустая сумма равна нулю: узел, не связанный путём с наблюдаемым
выходом, не вносит в него вклад. Для цепочки
$\mathcal{P}_l=\{l-1\}$ получаем произведение
$M_l=J^{(L,L-1)}\cdots J^{(l+1,l)}$ при $l<L$.
При ветвлении $M_l$ суммирует произведения по всем путям,
включая обходные связи.
Равенство $d=\Delta\Vec{y}^{(L)}$ для произвольной сети не предполагается.
Изолированно разность измеряется на опорном входе $\Vec{x}^{(l)}$;
на фактическом входе $\widetilde{\Vec{x}}^{(l)}$ она может измениться.

Для каждого фиксированного $M_l$ выполняется
$\|M_l\delta_l\|_2^2=\delta_l^TM_l^TM_l\delta_l$.
Поэтому одинаковая норма входной ошибки не определяет
одинаковое выходное расхождение: оно зависит и от направления ошибки.
Кроме направления отдельного вклада важна взаимная ориентация разных
вкладов: при сложении они могут усиливать или компенсировать друг друга.
Чтобы выразить эту зависимость в среднем, рассмотрим ансамбль~---
совокупность возможных возмущений с заданным распределением и конечными
вторыми моментами. Обозначим математическое ожидание по нему через
$\mathbb{E}$, дисперсию~--- через $\operatorname{Var}$ и положим
$\mu_l=\mathbb{E}\delta_l$,
$C_{lm}=\mathbb{E}[(\delta_l-\mu_l)(\delta_m-\mu_m)^T]$.
Это средние и перекрёстные ковариации локальных ошибок.
Пусть $\operatorname{tr}$ обозначает след матрицы, то есть сумму
её диагональных элементов; индексы $l,m$ далее пробегают $1,\ldots,L$.
При фиксированных $M_l$ тождества моментов~\cite[разд.~6.2]{Petersen2012}
дают средний квадрат нормы суммарного возмущения:
\begin{equation}
    \label{eq:moments}
    \mathbb{E}\|d\|_2^2=
    \left\|\sum_l M_l\mu_l\right\|_2^2
    +\sum_{l,m}\operatorname{tr}(M_l C_{lm}M_m^T),
\end{equation}
При усреднении по разным входам формула применяется условно
для каждого входа: матрицы переноса также зависят от него.

При равных сонаправленных перенесённых средних их сумма
в~\eqref{eq:moments} растёт как $L$; для некоррелированных центрированных
вкладов с одинаковой среднеквадратической нормой сумма растёт
как $\sqrt L$ в среднеквадратическом смысле.
Тот же условный вывод для суммы $K$ скалярных ошибок $e_k$
с общим средним $\mu_e$, дисперсией $\sigma_e^2$
и нулевыми попарными ковариациями имеет вид
\begin{equation}
    \label{eq:scaling}
    \mathbb{E}\sum_{k=1}^{K}e_k=K\mu_e,\qquad
    \sqrt{\operatorname{Var}\!\left(\sum_{k=1}^{K}e_k\right)}
    =\sqrt K\,\sigma_e.
\end{equation}
Зависимости~\eqref{eq:scaling} условны и не задают универсального закона относительной
ошибки GEMM.

Для экспериментальной проверки в опорное выполнение графа вводятся
измеренные локальные ошибки. Метрику~\eqref{eq:drel} итогового выхода
обозначим $d_{\mathrm{rel}}^{\text{сеть}}$. При ненулевом максимальном
локальном расхождении определим коэффициент $R_{\mathrm{fwd}}$:
\begin{equation}
    \label{eq:ratio}
    R_{\mathrm{fwd}}=
    \frac{d_{\mathrm{rel}}^{\text{сеть}}}
         {\max_l d_{\mathrm{rel}}^{(l)}},\qquad
    \max_l d_{\mathrm{rel}}^{(l)}>0,
\end{equation}
Коэффициент сопоставляет итоговую относительную ошибку с наибольшей
локальной. Их опорные нормы различаются, поэтому это не норма оператора
и $R_{\mathrm{fwd}}<1$ не доказывает наличие сжимающего слоя.

\subsection{Ослабление и взаимная компенсация ошибок}

Формула~\eqref{eq:accum} допускает два разных механизма уменьшения
итогового отклонения: последующие операции ослабляют отдельный вклад
или разные вклады компенсируются при сложении. Эти механизмы
не ограничиваются нормализацией.

\textbf{Ослабление отдельного возмущения.}
Рассмотрим операцию с $n\ge1$ входными и $m\ge1$ выходными координатами.
Пусть $u\in\mathbb{R}^{n}$~--- изменение входа,
$J\in\mathbb{R}^{m\times n}$~--- её якобиан в опорной точке.
Спектральная норма $\|J\|_2$ задаёт максимальное усиление евклидовой
нормы в линейном приближении:
\[
    \|Ju\|_2\le\|J\|_2\|u\|_2.
\]
При $\|J\|_2<1$ операция локально ослабляет любое направление.
Даже при большей норме конкретное $u$ может ослабляться
или полностью исчезать, если $Ju=0$.
Например, выборка координат удаляет компоненты ошибки,
относящиеся к отброшенным координатам.

Пусть $x_i$~--- координаты входа, $\phi'$~--- производная
поэлементной функции активации $\phi$. Тогда её якобиан диагонален
и имеет норму $\max_i|\phi'(x_i)|$.
Для сигмоиды $\sigma(t)=1/(1+\exp(-t))$ при вещественном $t$ производная равна
$\sigma(t)(1-\sigma(t))\le1/4$. Поэтому сама сигмоида ослабляет
абсолютное возмущение; предшествующее линейное преобразование при этом
может его усиливать. В общем случае ослабление активацией зависит
от её производной на текущем входе.

\textbf{Нормализация как частный случай.}
Нормализация слоя (Layer Normalization, LayerNorm)~\cite{Ba2016}
вычитает среднее по координатам. Нормализация по среднеквадратическому
масштабу (Root Mean Square Layer Normalization, RMSNorm)~\cite{Zhang2019}
не выполняет центрирование.
Для входа $\Vec{x}\in\mathbb{R}^{n}$ определим среднее
$\bar x=n^{-1}\sum_{i=1}^{n}x_i$ и вектор единиц $\Vec{1}\in\mathbb{R}^{n}$.
Центрирование можно записать как
\[
    z=\Vec{x}-\bar x\Vec{1}=P\Vec{x},\qquad
    P=I_n-\frac{\Vec{1}\Vec{1}^{T}}{n}.
\]
Матрица $P$ удаляет постоянную по координатам составляющую и сохраняет
векторы с нулевым средним. Такое преобразование называют ортогональной проекцией, а матрицу $P$~---
проектором; повторное центрирование ничего не меняет.
Для RMSNorm положим $P=I_n$, так что $z=\Vec{x}$.

Для предотвращения нулевого знаменателя введём $\epsilon_{\mathrm{norm}}>0$.
Оба слоя делят $z$ на входной масштаб
$s=(\|z\|_2^2/n+\epsilon_{\mathrm{norm}})^{1/2}$.
Пусть $\gamma,\beta\in\mathbb{R}^{n}$~--- фиксированные масштаб и сдвиг,
$D_\gamma=\operatorname{diag}(\gamma)$~--- диагональная матрица масштабов,
$\|\gamma\|_\infty=\max_i|\gamma_i|$.
Для выхода $\Vec{y}=D_\gamma z/s+\beta$ дифференцирование по входу даёт
\begin{equation}
    \label{eq:sink}
    J=\frac{D_\gamma}{s}
       \left(P-\frac{zz^T}{ns^2}\right),\qquad
    \|J\|_2\le\frac{\|\gamma\|_\infty}{s}.
\end{equation}
Член $P$ описывает выбранную обработку входа, а поправка
$zz^T/(ns^2)$~--- изменение самого делителя при возмущении входа.
Поскольку $z$ лежит в сохраняемом $P$ подпространстве и
$\|z\|_2^2/(ns^2)\le1$, норма скобки не превышает единицы.
Из оценки~\eqref{eq:sink} следует достаточное условие.
\begin{statement}[Достаточное условие сжатия]\label{st:sink}
    При $s>\|\gamma\|_\infty$ нормализационный слой локально сжимает
    евклидову норму возмущения.
\end{statement}
Это условие зависит от входа. Кроме того, остаточный блок
$\Vec{x}+F(\Vec{x})$, где $F$~--- преобразование ветви,
имеет якобиан $I_n+J_F$, где $J_F$~--- якобиан $F$.
Поэтому ослабление внутри ветви не гарантирует ослабления всего блока.

\textbf{Компенсация разных источников.}
Пусть $v_l=M_l\delta_l$~--- вклад слоя $l$, уже перенесённый к выходу.
Обозначим скалярное произведение через $\langle v_l,v_m\rangle=v_l^Tv_m$.
Разложение квадрата нормы суммы имеет вид
\[
    \|d\|_2^2
    =\sum_l\|v_l\|_2^2
     +2\sum_{l<m}\langle v_l,v_m\rangle.
\]
Отрицательные перекрёстные слагаемые уменьшают итоговую норму
по сравнению с суммой квадратов отдельных вкладов.
Такая компенсация возможна даже без сжимающего слоя.
Она зависит от направлений и взаимной зависимости ошибок,
которые входят в~\eqref{eq:moments}; локальных норм для её предсказания
недостаточно. Для полной нелинейной сети отдельно остаётся проверить
остаток $r_{\mathrm{lin}}$.

\subsection{Распространение ошибок в обратном проходе}

В обратном проходе градиенты вычисляются по сохранённым тензорам
прямого прохода, которые также могут содержать ошибки.
Сначала зафиксируем опорные тензоры.

Пусть $\mathcal{L}$~--- скалярная функция потерь от наблюдаемого выхода,
$\Vec{g}^{(j)}=\nabla_{\Vec{y}^{(j)}}\mathcal{L}\in\mathbb{R}^{n_j}$
при $j=0,\ldots,L$.
При обратном проходе узел получает градиентные вклады от всех
потребителей своего выхода, поэтому они суммируются:
\[
    \Vec{g}^{(j)}=\sum_{k\in\mathcal{S}_j}J^{(k,j)T}\Vec{g}^{(k)},
    \qquad 0\le j<L.
\]
Для $p_l>0$ обозначим якобиан по параметрам через
$H^{(l)}=\partial f^{(l)}/\partial\Vec{w}^{(l)}
\in\mathbb{R}^{n_l\times p_l}$,
а градиент по ним~--- через
$\Vec{q}^{(l)}=H^{(l)T}\Vec{g}^{(l)}\in\mathbb{R}^{p_l}$.
При общем параметре нескольких узлов их градиентные вклады складываются.
Пусть $h^{(j)}=\widetilde{\Vec{g}}^{(j)}-\Vec{g}^{(j)}$~--- отклонение
передаваемого градиента,
$\eta_x^{(j)}\in\mathbb{R}^{n_j}$~--- собственная ошибка вычисления
и суммирования вкладов в узле $j$,
$\eta_w^{(l)}\in\mathbb{R}^{p_l}$~--- ошибка вычисления градиента
по параметрам. На фиксированных тензорах прямого прохода
\begin{equation}
    \label{eq:bwd}
    \begin{aligned}
        h^{(j)}&=\sum_{k\in\mathcal{S}_j}J^{(k,j)T}h^{(k)}
            +\eta_x^{(j)},\qquad 0\le j<L,\\
        \widetilde{\Vec{q}}^{(l)}-\Vec{q}^{(l)}
            &=H^{(l)T}h^{(l)}+\eta_w^{(l)},\qquad p_l>0.
    \end{aligned}
\end{equation}
Для графа на рис.~\ref{fig:transport} обратное распространение
погрешностей градиентов по активациям показано на рис.~\ref{fig:backward-dag}.

\begingroup
\setlength{\intextsep}{3pt}
\begin{figure}[H]
    \centering
    \begin{tikzpicture}[
        font=\normalsize,
        box/.style={draw=black!65, rounded corners=2pt,
            fill=babyblue!20, minimum width=16mm, minimum height=8mm},
        arr/.style={-{Stealth[length=4pt]}, line width=0.6pt},
        loc/.style={arr, red!70!black}
    ]
        \node (g0) at (-2.6,0) {$h^{(0)}$};
        \node[box] (g1) at (0,0) {$h^{(1)}$};
        \node[box] (g2) at (3,0.9) {$h^{(2)}$};
        \node[box] (g3) at (3,-0.9) {$h^{(3)}$};
        \node[box] (g4) at (6,0) {$h^{(4)}$};
        \draw[arr] (g1.west) -- node[above] {$J^{(1,0)T}$} (g0.east);
        \draw[arr] (g2.west) -- node[above, sloped] {$J^{(2,1)T}$} (g1.east);
        \draw[arr] (g3.west) -- node[below, sloped] {$J^{(3,1)T}$} (g1.east);
        \draw[arr] (g4.west) -- node[above, sloped] {$J^{(4,2)T}$} (g2.east);
        \draw[arr] (g4.west) -- node[below, sloped] {$J^{(4,3)T}$} (g3.east);
        \node (e0) at (-2.6,1.1) {$\eta_x^{(0)}$};
        \node (e1) at (0,1.1) {$\eta_x^{(1)}$};
        \node (e2) at (4.4,1.4) {$\eta_x^{(2)}$};
        \node (e3) at (4.4,-1.4) {$\eta_x^{(3)}$};
        \draw[loc] (e0) -- (g0.north);
        \draw[loc] (e1) -- (g1.north);
        \draw[loc] (e2) -- (g2.north);
        \draw[loc] (e3) -- (g3.south);
        \node at (2,-2.2) {
            $h^{(1)}=J^{(2,1)T}h^{(2)}
                +J^{(3,1)T}h^{(3)}+\eta_x^{(1)}$
        };
    \end{tikzpicture}
    \caption{Обратное распространение погрешностей градиентов для графа
        на рис.~\ref{fig:transport}. На рёбрах указаны транспонированные
        якобианы; красные стрелки задают локальные ошибки $\eta_x^{(j)}$.
        Начальное $h^{(4)}$ включает ошибку градиента функции потерь.}
    \label{fig:backward-dag}
\end{figure}
\endgroup

Начальное $h^{(L)}$ включает возможное расхождение градиента функции потерь.
При изменённых тензорах прямого прохода меняются также якобианы.

Чтобы отделить собственные ошибки обратных вычислений от влияния
изменённых тензоров прямого прохода, нужны два режима сравнения.
Пусть $q$~--- один выбранный градиентный тензор,
$\widetilde q_{\mathrm{iso}}$~--- его целевое значение на опорных тензорах
прямого прохода, $\widetilde q_{\mathrm{joint}}$~--- на собственных тензорах.
Обозначим собственное отклонение обратных вычислений через
$\Delta_{\mathrm{bwd}}$, а дополнительное изменение при замене тензоров
прямого прохода~--- через $\Delta_{\mathrm{fwd}}$.
Прибавление и вычитание $\widetilde q_{\mathrm{iso}}$ даёт точное разложение:
\begin{equation}
    \label{eq:grad-decomp}
    \widetilde q_{\mathrm{joint}}-q
    =\underbrace{(\widetilde q_{\mathrm{iso}}-q)}_{\Delta_{\mathrm{bwd}}}
    +\underbrace{(\widetilde q_{\mathrm{joint}}-\widetilde q_{\mathrm{iso}})}
        _{\Delta_{\mathrm{fwd}}}.
\end{equation}
Вклады могут усиливать или компенсировать друг друга;
для их норм выполняется лишь неравенство
$\|\Delta_{\mathrm{bwd}}+\Delta_{\mathrm{fwd}}\|_F
\le\|\Delta_{\mathrm{bwd}}\|_F+\|\Delta_{\mathrm{fwd}}\|_F$,
которое не предписывает, в каком режиме расхождение больше.

\subsection{Диагностические характеристики}

Норма в~\eqref{eq:drel} показывает величину расхождения.
Однако реализации с близкими нормами могут различаться по знакам
отклонений и их связи с опорными значениями. Поэтому для описания
характера ошибок дополнительно используются два показателя:
доля положительных отклонений $b$ и выборочная корреляция Пирсона $c$,
характеризующая линейную связь ошибки с опорой.

Для $n\ge1$ элементов сравниваемого тензора обозначим опорные
и сравниваемые значения через $y_i$ и $\widetilde y_i$,
их разности~--- через $e_i=\widetilde y_i-y_i$, $i=1,\ldots,n$.
Доля $b$ рассчитывается по всем $n$ элементам.
Индикатор $\mathbf{1}$ равен единице при выполнении условия и нулю иначе:
\begin{equation}
    \label{eq:diag}
    b=\frac1n\sum_{i=1}^{n}\mathbf{1}[e_i>0],\qquad
    c=\operatorname{corr}(e_i,y_i).
\end{equation}
Корреляция определена при $n\ge2$ и ненулевых дисперсиях обоих наборов.
Доля $b$ не измеряет среднее смещение ошибки, а $c\approx0$ не доказывает
независимости. Эти показатели описывают расхождения в одном тензоре;
предположения об ансамбле в~\eqref{eq:moments} требуют отдельной проверки.
